\section{Published figures of other systems}
\label{sec:published}

Table~\ref{tab:published} lists \nPublishedRows{} figures published by other work on these
benchmarks: \nPublishedHotpotqa{} on HotpotQA, \nPublishedMusique{} on MuSiQue and \nPublishedMultihop{}
on MultiHop-RAG. They cover the BM25-plus-cross-encoder recipe and dense baselines of
BEIR~\citep{thakur2021beir}, trained multi-hop retrievers
(MDR~\citep{xiong2021mdr}, Baleen~\citep{khattab2021baleen}), retrieval interleaved with
language-model reasoning (IRCoT~\citep{trivedi2023ircot}, HippoRAG~\citep{gutierrez2024hipporag,gutierrez2025hipporag2})
and the MultiHop-RAG paper's own pipelines~\citep{tang2024multihoprag}. Each row was read in the
paper it comes from; \nPublishedExcluded{} candidate figures that could not be verified there were
excluded rather than quoted from a survey or another paper's table. The appendix repeats every row with its
corpus and a note on how it differs from our setting (Tables~\ref{tab:published-hotpotqa}
to~\ref{tab:published-multihop}).

\begin{table*}[!t]
\centering
\scriptsize
\setlength{\tabcolsep}{3pt}
\caption{\capPublishedShort}
\label{tab:published}
% Generated by `cer p17-tables`; do not edit.
% caption: Published figures of other systems on the three benchmarks, not a paired comparison with ours. Each row was read in the paper itself; the metric and the unit are the paper's own, and how each differs from our setting is in the appendix.
\begin{tabular}{p{5.2cm}p{6.4cm}rrp{2.1cm}l}
\toprule
System & Metric & Value & Questions & Training or LLM & Source \\
\midrule
\multicolumn{6}{l}{\textbf{HotpotQA}} \\
BM25 & nDCG@10 & 0.603 & 7,405 & none & \cite{thakur2021beir} \\
BM25 + cross-encoder reranking (MiniLM-L6, top-100) & nDCG@10 & 0.707 & 7,405 & trained retriever & \cite{thakur2021beir} \\
DPR & nDCG@10 & 0.391 & 7,405 & trained retriever & \cite{thakur2021beir} \\
ANCE & nDCG@10 & 0.456 & 7,405 & trained retriever & \cite{thakur2021beir} \\
TAS-B & nDCG@10 & 0.584 & 7,405 & trained retriever & \cite{thakur2021beir} \\
ColBERT (late interaction, end-to-end) & nDCG@10 & 0.593 & 7,405 & trained retriever & \cite{thakur2021beir} \\
MDR (multi-hop dense retriever, trained on HotpotQA) & R@2 (recall at k retrieved passages, percent) & 65.9 & -- & trained retriever & \cite{xiong2021mdr} \\
MDR (multi-hop dense retriever, trained on HotpotQA) & R@10 (recall at k retrieved passages, percent) & 77.5 & -- & trained retriever & \cite{xiong2021mdr} \\
MDR (multi-hop dense retriever, trained on HotpotQA) & R@20 (recall at k retrieved passages, percent) & 80.2 & -- & trained retriever & \cite{xiong2021mdr} \\
Baleen (FLIPR + condenser, trained on HotpotQA) & P-EM (both gold passages exactly selected, percent) & 86.7 & 7,405 & trained retriever & \cite{khattab2021baleen} \\
Baleen (FLIPR + condenser, trained on HotpotQA) & P-R@20 (top-20 passages contain both gold passages, percent) & 93.3 & 7,405 & trained retriever & \cite{khattab2021baleen} \\
BM25 & passage recall@2 (percent) on HotpotQA & 55.4 & 1,000 & none & \cite{gutierrez2024hipporag} \\
BM25 & passage recall@5 (percent) on HotpotQA & 72.2 & 1,000 & none & \cite{gutierrez2024hipporag} \\
ColBERTv2 & passage recall@2 (percent) on HotpotQA & 64.7 & 1,000 & trained retriever & \cite{gutierrez2024hipporag} \\
ColBERTv2 & passage recall@5 (percent) on HotpotQA & 79.3 & 1,000 & trained retriever & \cite{gutierrez2024hipporag} \\
HippoRAG (ColBERTv2) & passage recall@2 (percent) on HotpotQA & 60.5 & 1,000 & LLM in the loop & \cite{gutierrez2024hipporag} \\
HippoRAG (ColBERTv2) & passage recall@5 (percent) on HotpotQA & 77.7 & 1,000 & LLM in the loop & \cite{gutierrez2024hipporag} \\
IRCoT + ColBERTv2 & passage recall@2 (percent) on HotpotQA & 67.9 & 1,000 & LLM in the loop & \cite{gutierrez2024hipporag} \\
IRCoT + ColBERTv2 & passage recall@5 (percent) on HotpotQA & 82 & 1,000 & LLM in the loop & \cite{gutierrez2024hipporag} \\
IRCoT + HippoRAG (ColBERTv2) & passage recall@2 (percent) on HotpotQA & 67 & 1,000 & LLM in the loop & \cite{gutierrez2024hipporag} \\
IRCoT + HippoRAG (ColBERTv2) & passage recall@5 (percent) on HotpotQA & 83 & 1,000 & LLM in the loop & \cite{gutierrez2024hipporag} \\
BM25 & passage recall@2 (percent) on HotpotQA & 57.3 & 1,000 & none & \cite{gutierrez2025hipporag2} \\
BM25 & passage recall@5 (percent) on HotpotQA & 74.8 & 1,000 & none & \cite{gutierrez2025hipporag2} \\
NV-Embed-v2 (7B dense retriever) & passage recall@2 (percent) on HotpotQA & 84.1 & 1,000 & trained retriever & \cite{gutierrez2025hipporag2} \\
NV-Embed-v2 (7B dense retriever) & passage recall@5 (percent) on HotpotQA & 94.5 & 1,000 & trained retriever & \cite{gutierrez2025hipporag2} \\
HippoRAG 2 (Llama-3.3-70B-Instruct) & passage recall@2 (percent) on HotpotQA & 83.5 & 1,000 & LLM in the loop & \cite{gutierrez2025hipporag2} \\
HippoRAG 2 (Llama-3.3-70B-Instruct) & passage recall@5 (percent) on HotpotQA & 96.3 & 1,000 & LLM in the loop & \cite{gutierrez2025hipporag2} \\
\midrule
\multicolumn{6}{l}{\textbf{MuSiQue}} \\
BM25 & passage recall@2 (percent) on MuSiQue (answerable) & 32.3 & 1,000 & none & \cite{gutierrez2024hipporag} \\
BM25 & passage recall@5 (percent) on MuSiQue (answerable) & 41.2 & 1,000 & none & \cite{gutierrez2024hipporag} \\
ColBERTv2 & passage recall@2 (percent) on MuSiQue (answerable) & 37.9 & 1,000 & trained retriever & \cite{gutierrez2024hipporag} \\
ColBERTv2 & passage recall@5 (percent) on MuSiQue (answerable) & 49.2 & 1,000 & trained retriever & \cite{gutierrez2024hipporag} \\
HippoRAG (ColBERTv2) & passage recall@2 (percent) on MuSiQue (answerable) & 40.9 & 1,000 & LLM in the loop & \cite{gutierrez2024hipporag} \\
HippoRAG (ColBERTv2) & passage recall@5 (percent) on MuSiQue (answerable) & 51.9 & 1,000 & LLM in the loop & \cite{gutierrez2024hipporag} \\
IRCoT + ColBERTv2 & passage recall@2 (percent) on MuSiQue (answerable) & 41.7 & 1,000 & LLM in the loop & \cite{gutierrez2024hipporag} \\
IRCoT + ColBERTv2 & passage recall@5 (percent) on MuSiQue (answerable) & 53.7 & 1,000 & LLM in the loop & \cite{gutierrez2024hipporag} \\
IRCoT + HippoRAG (ColBERTv2) & passage recall@2 (percent) on MuSiQue (answerable) & 45.3 & 1,000 & LLM in the loop & \cite{gutierrez2024hipporag} \\
IRCoT + HippoRAG (ColBERTv2) & passage recall@5 (percent) on MuSiQue (answerable) & 57.6 & 1,000 & LLM in the loop & \cite{gutierrez2024hipporag} \\
BM25 & passage recall@2 (percent) on MuSiQue & 32.4 & 1,000 & none & \cite{gutierrez2025hipporag2} \\
BM25 & passage recall@5 (percent) on MuSiQue & 43.5 & 1,000 & none & \cite{gutierrez2025hipporag2} \\
NV-Embed-v2 (7B dense retriever) & passage recall@2 (percent) on MuSiQue & 52.7 & 1,000 & trained retriever & \cite{gutierrez2025hipporag2} \\
NV-Embed-v2 (7B dense retriever) & passage recall@5 (percent) on MuSiQue & 69.7 & 1,000 & trained retriever & \cite{gutierrez2025hipporag2} \\
HippoRAG 2 (Llama-3.3-70B-Instruct) & passage recall@2 (percent) on MuSiQue & 56.1 & 1,000 & LLM in the loop & \cite{gutierrez2025hipporag2} \\
HippoRAG 2 (Llama-3.3-70B-Instruct) & passage recall@5 (percent) on MuSiQue & 74.7 & 1,000 & LLM in the loop & \cite{gutierrez2025hipporag2} \\
\midrule
\multicolumn{6}{l}{\textbf{MultiHop-RAG}} \\
text-embedding-ada-002 (without reranker) & Hits@10 & 0.6381 & 2,255 & trained retriever & \cite{tang2024multihoprag} \\
text-embedding-ada-002 (without reranker) & Hits@4 & 0.504 & 2,255 & trained retriever & \cite{tang2024multihoprag} \\
text-embedding-ada-002 (with bge-reranker-large) & Hits@10 & 0.7059 & 2,255 & trained retriever & \cite{tang2024multihoprag} \\
text-embedding-ada-002 (with bge-reranker-large) & Hits@4 & 0.6169 & 2,255 & trained retriever & \cite{tang2024multihoprag} \\
bge-large-en-v1.5 (without reranker) & Hits@10 & 0.6718 & 2,255 & trained retriever & \cite{tang2024multihoprag} \\
bge-large-en-v1.5 (without reranker) & Hits@4 & 0.5221 & 2,255 & trained retriever & \cite{tang2024multihoprag} \\
bge-large-en-v1.5 (with bge-reranker-large) & Hits@10 & 0.7183 & 2,255 & trained retriever & \cite{tang2024multihoprag} \\
bge-large-en-v1.5 (with bge-reranker-large) & Hits@4 & 0.6364 & 2,255 & trained retriever & \cite{tang2024multihoprag} \\
voyage-02 (without reranker) & Hits@10 & 0.6506 & 2,255 & trained retriever & \cite{tang2024multihoprag} \\
voyage-02 (without reranker) & Hits@4 & 0.4619 & 2,255 & trained retriever & \cite{tang2024multihoprag} \\
voyage-02 (with bge-reranker-large) & Hits@10 & 0.7467 & 2,255 & trained retriever & \cite{tang2024multihoprag} \\
voyage-02 (with bge-reranker-large) & Hits@4 & 0.6625 & 2,255 & trained retriever & \cite{tang2024multihoprag} \\
\bottomrule
\end{tabular}

\end{table*}

\paragraph{What the reader may conclude.} Our nDCG@10 and gold-recall figures
(Appendix~\ref{app:d6}) use the corpus unit, the whole paragraph, and the same questions
and corpus as BEIR on HotpotQA, so they can be laid beside the BM25-plus-cross-encoder row as
context: for instance, our J-light over Dense plus BM25 reaches \ndcgLightPTenBHotpot{} (nDCG@10, times
one hundred), and J-strong \ndcgStrongPTenBHotpot{}. They place the zero-shot recipe of this paper in the
same range as other zero-shot recipes on identical questions.

\paragraph{What the reader may not conclude.} These are \emph{published figures, not a paired
comparison}. The rows differ from our setting in at least one of the unit of retrieval
(chunks of a fixed token length, or whole paragraphs), the metric (nDCG with graded relevance, recall at
a rank, hits at a rank, or our token-budget Full Support), the corpus (a pooled
corpus of a subset of questions, or the whole of Wikipedia), the questions (a sampled subset, or all of
them), and whether the system is trained or calls a language model; the differences column says which,
row by row. Our weights are fitted on HotpotQA development questions, so our HotpotQA figures are
in-sample for P10-C and P14. Our systems are untrained, call no model at query time, and are
cheap; they are not offered as competitors of trained or language-model-based systems, and
nothing in this paper shows that they match or beat those systems.
